\documentclass[11pt]{article}
\usepackage[a4paper,margin=22mm]{geometry}
\usepackage{fontspec}
\setmainfont{DejaVu Serif}
\setsansfont{DejaVu Sans}
\usepackage{microtype}
\usepackage{xcolor}
\usepackage{hyperref}
\usepackage{enumitem}
\usepackage{parskip}
\definecolor{Accent}{HTML}{971D32}
\definecolor{Muted}{HTML}{666666}
\hypersetup{colorlinks=true,urlcolor=Accent,linkcolor=Accent}
\setlist{leftmargin=1.6em,itemsep=0.3em,topsep=0.35em}
\setcounter{tocdepth}{2}
\renewcommand{\contentsname}{Contents}
\newcommand{\sectionrule}{\par\vspace{-0.5em}\noindent\textcolor{Accent}{\rule{\linewidth}{0.6pt}}\par}
\begin{document}

\begin{center}
{\sffamily\bfseries\color{Accent} TOEFL WORD OF THE DAY\par}
\vspace{8mm}
{\fontsize{42}{48}\selectfont\bfseries misattribute\par}
\vspace{2mm}
{\Large\sffamily\color{Accent} /ˌmɪsəˈtrɪbjuːt/\par}
\vspace{2mm}
{\small\sffamily Local audio phrase: misattribute\par}
{\small\sffamily Audio file: audios/misattribute.mp3\par}
\vspace{2mm}
{\large\sffamily\color{Muted} transitive verb\par}
\vspace{5mm}
{\sffamily identify the wrong cause \textbullet\ identify the wrong creator or source\par}
\vspace{6mm}
{\large To misattribute something is to connect it with the wrong cause, creator, speaker, or source.\par}
\vspace{6mm}
\begin{minipage}{0.84\textwidth}
\itshape “The quotation is frequently misattributed to Albert Einstein.”
\end{minipage}\par
\vspace{10mm}
\textbf{Date:} August 13th 2026\quad\textbullet\quad \textbf{Level:} B1--mid-B2 TOEFL
\vfill
Created and compiled by \textbf{Amir Kooshky}\par
\vspace{2mm}
\href{https://t.me/KooshkyTOEFL}{Telegram Channel}\quad\textbullet\quad
\href{https://www.instagram.com/kooshkytoefl}{Instagram}\quad\textbullet\quad
\href{https://t.me/KooshkyTOEFL\_pv}{Personal Telegram}
\end{center}
\newpage
\tableofcontents
\newpage

\section{Meanings and usage}
\sectionrule
\subsection{Meaning 1: Identify the wrong cause}
\textbf{Part of speech:} transitive verb

\textbf{IPA:} /ˌmɪsəˈtrɪbjuːt/

\textbf{Pronunciation phrase:} misattribute

\textbf{Local audio file:} audios/misattribute.mp3

\textbf{Register:} formal; common in academic, scientific, medical, psychological, and analytical writing

\textbf{Definition:} to wrongly believe that an event, condition, or result was caused by a particular person, thing, or factor

\textbf{In simpler words:} blame or credit the wrong cause

\textbf{Typical pattern:} misattribute + result, change, problem, or effect + to + wrong cause

\subsubsection{Natural examples}
\begin{enumerate}
  \item Researchers initially misattributed the population decline to disease rather than habitat loss.
  \item People sometimes misattribute normal changes in the weather to long-term climate change.
  \item The company misattributed the drop in sales to higher prices, but later research showed that poor customer service was the main cause.
  \item Doctors may misattribute physical symptoms to stress if they do not consider other possible causes.
  \item The researchers warned against misattributing differences between the groups to a single factor.
  \item Observers misattributed the improvement to the new policy even though the trend had begun several years earlier.
  \item Students often misattribute poor test performance to lack of ability when ineffective study methods may be responsible.
  \item The early study misattributed the increase in mortality to pollution because an important variable had been overlooked.
\end{enumerate}

\subsubsection{Nuance and implication}
\textbf{Central nuance:} The word means that a connection between a result and its cause has been made incorrectly.

\textbf{Usual implication:} The mistake often happens because evidence is incomplete, another factor has been overlooked, or correlation has been confused with causation.

\textbf{Compared with attribute:} To attribute a result to something is to identify it as the believed cause. To misattribute it is to identify the wrong cause.

\subsubsection{Grammar notes}
\textbf{How to use it:} Place the result, event, or condition directly after misattribute.

\textbf{What can follow it:} After the result, use to and then the cause that has been identified incorrectly.

\textbf{Common preposition:} Use misattribute something to something, as in misattribute the decline to pollution.

\textbf{Position in a sentence:} The passive form is also common: The decline was misattributed to pollution.

\subsubsection{Useful patterns}
\textbf{Pattern:} misattribute + result + to + cause

\textbf{Use:} Use this when someone incorrectly identifies what caused a result.

\textbf{Example:} Analysts misattributed the economic slowdown to lower consumer demand.

\textbf{Pattern:} be misattributed to + cause

\textbf{Use:} Use the passive form when the incorrect explanation is more important than the person making it.

\textbf{Example:} The symptoms were initially misattributed to anxiety.

\textbf{Pattern:} misattribute a difference to + factor

\textbf{Use:} Use this when someone incorrectly explains why two groups, periods, or conditions differ.

\textbf{Example:} The researchers cautioned against misattributing the difference to age alone.

\textbf{Pattern:} wrongly or mistakenly misattribute

\textbf{Use:} Usually avoid this pattern because misattribute already contains the idea of being wrong.

\textbf{Example:} The study misattributed the effect to changes in temperature.

\subsubsection{Common wrong usage}
\paragraph{Correction 1}
\textbf{Wrong:} The researchers misattributed the decline from pollution.

\textbf{Correct:} The researchers misattributed the decline to pollution.

\textbf{Why:} Use misattribute something to a cause, not from.

\paragraph{Correction 2}
\textbf{Wrong:} The researchers misattributed that pollution caused the decline.

\textbf{Correct:} The researchers misattributed the decline to pollution.

\textbf{Why:} The usual pattern is misattribute the result to the supposed cause.

\paragraph{Correction 3}
\textbf{Wrong:} The decline was misattributed by climate change.

\textbf{Correct:} The decline was misattributed to climate change.

\textbf{Why:} Use to before the incorrectly identified cause. By would introduce the person making the attribution.

\paragraph{Correction 4}
\textbf{Wrong:} Scientists wrongly misattributed the effect to diet.

\textbf{Correct:} Scientists misattributed the effect to diet.

\textbf{Why:} Misattribute already means attribute incorrectly, so wrongly is usually unnecessary.

\subsection{Meaning 2: Identify the wrong creator or source}
\textbf{Part of speech:} transitive verb

\textbf{IPA:} /ˌmɪsəˈtrɪbjuːt/

\textbf{Pronunciation phrase:} misattribute

\textbf{Local audio file:} audios/misattribute.mp3

\textbf{Register:} formal; common in history, literature, art, journalism, scholarship, and discussions of sources

\textbf{Definition:} to wrongly identify a particular person or group as the creator, author, speaker, or source of a work, idea, statement, or achievement

\textbf{In simpler words:} say the wrong person created or said something

\textbf{Typical pattern:} misattribute + quotation, work, idea, discovery, or statement + to + person

\subsubsection{Natural examples}
\begin{enumerate}
  \item The quotation is frequently misattributed to Albert Einstein.
  \item For many years, scholars misattributed the painting to a more famous artist.
  \item Historical inventions are sometimes misattributed to a single individual even when many people contributed to their development.
  \item The poem was misattributed to a nineteenth-century writer because of a cataloging error.
  \item Several online articles misattribute the statement to a former president.
  \item Researchers discovered that the manuscript had been misattributed to the wrong author.
  \item The discovery is often misattributed to one scientist, although two independent teams reported it at nearly the same time.
  \item A widely shared social media post misattributed the photograph to a completely different event.
\end{enumerate}

\subsubsection{Nuance and implication}
\textbf{Central nuance:} The word means that the source, creator, speaker, or person responsible for something has been identified incorrectly.

\textbf{Usual implication:} A misattribution may result from weak evidence, repetition of an earlier mistake, incomplete records, or deliberate misinformation.

\textbf{Compared with misquote:} To misquote someone is to report their words inaccurately. To misattribute a quotation is to identify the wrong person as the one who said or wrote it.

\subsubsection{Grammar notes}
\textbf{How to use it:} Place the work, statement, idea, discovery, or other item directly after misattribute.

\textbf{What can follow it:} Use to and then the person, group, culture, or source that has been incorrectly identified.

\textbf{Common preposition:} Use misattribute something to someone, as in misattribute the quotation to Einstein.

\textbf{Position in a sentence:} The passive form is extremely common, especially with quotations and works: The quotation is often misattributed to Einstein.

\subsubsection{Useful patterns}
\textbf{Pattern:} misattribute + quotation or statement + to + person

\textbf{Use:} Use this when the wrong person is identified as having said or written something.

\textbf{Example:} Many websites misattribute the quotation to Mark Twain.

\textbf{Pattern:} misattribute + work + to + artist or author

\textbf{Use:} Use this when the wrong person is identified as the creator of a work.

\textbf{Example:} Early scholars misattributed the manuscript to a medieval poet.

\textbf{Pattern:} be commonly or frequently misattributed to + person

\textbf{Use:} Use this when the same incorrect attribution is widespread.

\textbf{Example:} The saying is commonly misattributed to Benjamin Franklin.

\textbf{Pattern:} be misattributed because of + reason

\textbf{Use:} Use this to explain why the source was identified incorrectly.

\textbf{Example:} The painting was misattributed because of an incorrect museum record.

\subsubsection{Common wrong usage}
\paragraph{Correction 1}
\textbf{Wrong:} The quotation is misattributed from Einstein.

\textbf{Correct:} The quotation is misattributed to Einstein.

\textbf{Why:} Use to before the incorrectly identified speaker, writer, or creator.

\paragraph{Correction 2}
\textbf{Wrong:} The historian misattributed the author for the poem.

\textbf{Correct:} The historian misattributed the poem to the author.

\textbf{Why:} Place the work directly after misattribute and the supposed creator after to.

\paragraph{Correction 3}
\textbf{Wrong:} The quotation was misquoted to Einstein.

\textbf{Correct:} The quotation was misattributed to Einstein.

\textbf{Why:} Use misattribute when the problem is the identity of the speaker. Misquote means report someone's words incorrectly.

\paragraph{Correction 4}
\textbf{Wrong:} The painting was wrongly misattributed to Picasso.

\textbf{Correct:} The painting was misattributed to Picasso.

\textbf{Why:} Misattributed already means attributed incorrectly, so wrongly is normally redundant.

\section{Word family}
\sectionrule
\subsection{misattribution}
\textbf{Part of speech:} countable or uncountable noun

\textbf{IPA:} /ˌmɪsˌætrəˈbjuːʃən/

\textbf{Definition:} an incorrect explanation of the cause of something or an incorrect identification of its creator, speaker, or source

\textbf{Relationship to the headword:} This noun names the mistake made when something is misattributed.

\textbf{Grammar pattern:} the misattribution of + result or work + to + cause or person

\textbf{Usage note:} It can be countable for one specific incorrect attribution and uncountable when discussing the phenomenon generally.

\subsubsection{Examples}
\begin{enumerate}
  \item The misattribution of the disease to poor diet delayed the search for its actual cause.
  \item The museum corrected the misattribution of the painting after new evidence identified its true artist.
  \item Online quotation databases contain many cases of misattribution.
\end{enumerate}

\section{Common collocations}
\sectionrule
\subsection{misattribute a cause}
\textbf{Applies to:} Identify the wrong cause

\textbf{Meaning:} identify the wrong explanation for an event or result

\textbf{Example:} People may misattribute a cause when several factors change at the same time.

\subsection{misattribute a change to}
\textbf{Applies to:} Identify the wrong cause

\textbf{Meaning:} incorrectly identify what caused a change

\textbf{Pattern:} misattribute a change to + supposed cause

\textbf{Example:} Analysts initially misattributed the change in consumer behavior to higher prices.

\subsection{misattribute symptoms to}
\textbf{Applies to:} Identify the wrong cause

\textbf{Meaning:} incorrectly identify the cause of symptoms

\textbf{Pattern:} misattribute symptoms to + condition or cause

\textbf{Example:} Doctors must be careful not to misattribute physical symptoms to stress.

\subsection{be misattributed to}
\textbf{Applies to:} all senses

\textbf{Meaning:} be incorrectly connected with a particular cause, creator, or source

\textbf{Pattern:} be misattributed to + cause, person, or source

\textbf{Example:} The improvement was initially misattributed to the new treatment.

\subsection{commonly misattributed to}
\textbf{Applies to:} Identify the wrong creator or source

\textbf{Meaning:} be widely but incorrectly believed to come from a particular person

\textbf{Pattern:} be commonly misattributed to + person

\textbf{Example:} The quotation is commonly misattributed to a famous philosopher.

\subsection{frequently misattributed to}
\textbf{Applies to:} Identify the wrong creator or source

\textbf{Meaning:} be repeatedly attributed to the wrong source

\textbf{Example:} Several ideas are frequently misattributed to the scientist who later popularized them.

\subsection{misattribute a quotation}
\textbf{Applies to:} Identify the wrong creator or source

\textbf{Meaning:} identify the wrong speaker or writer of a quotation

\textbf{Pattern:} misattribute a quotation to + person

\textbf{Example:} Writers should verify their sources so that they do not misattribute a quotation.

\subsection{misattribute a work}
\textbf{Applies to:} Identify the wrong creator or source

\textbf{Meaning:} identify the wrong creator of a work

\textbf{Pattern:} misattribute a work to + creator

\textbf{Example:} A damaged signature caused experts to misattribute the work to another painter.

\subsection{a case of misattribution}
\textbf{Applies to:} Identify the wrong creator or source

\textbf{Meaning:} an example in which the wrong source or creator has been identified

\textbf{Example:} The supposed quotation turned out to be a case of misattribution.

\textbf{Usage note:} Misattribution is the noun form.

\section{Synonyms and antonyms}
\sectionrule
\subsection{Close synonyms}
\subsubsection{mistakenly ascribe}
\textbf{Part of speech:} verb phrase

\textbf{Applies to:} all senses

\textbf{Shared meaning:} Both expressions mean incorrectly connect something with a cause, creator, or source.

\textbf{Difference:} Mistakenly ascribe has essentially the same meaning but is more formal and less compact. Misattribute directly expresses that the attribution is incorrect.

\textbf{Example:} Early researchers mistakenly ascribed the effect to changes in temperature.

\subsubsection{misidentify}
\textbf{Part of speech:} transitive verb

\textbf{Applies to:} Identify the wrong creator or source

\textbf{Shared meaning:} Both involve making an incorrect identification.

\textbf{Difference:} Misidentify broadly means identify any person or thing incorrectly. Misattribute specifically means connect a work, statement, result, or idea with the wrong source or creator.

\textbf{Example:} Witnesses initially misidentified the person in the photograph.

\section{Words learners may confuse}
\sectionrule
\subsection{misquote}
\textbf{Part of speech:} transitive verb

\textbf{Why learners confuse them:} Both commonly occur when discussing inaccurate quotations.

\textbf{Key difference:} Misattribute means identify the wrong person as the source of a quotation. Misquote means give the person's actual words inaccurately.

\textbf{misattribute:} The quotation is often misattributed to Einstein.

\textbf{misquote:} The article misquoted the scientist by changing several important words.

\subsection{misinterpret}
\textbf{Part of speech:} transitive verb

\textbf{Why learners confuse them:} Both describe errors in understanding evidence, but the type of error is different.

\textbf{Key difference:} Misattribute means connect something with the wrong cause or source. Misinterpret means understand the meaning or significance of something incorrectly.

\textbf{misattribute:} The researchers misattributed the decline to disease.

\textbf{misinterpret:} The researchers misinterpreted the decline as evidence that the species was disappearing.

\subsection{attribute}
\textbf{Part of speech:} transitive verb

\textbf{Why learners confuse them:} Misattribute is formed from attribute with the prefix mis-, meaning wrongly or incorrectly.

\textbf{Key difference:} Attribute means connect something with a believed cause or source. Misattribute means that this connection is incorrect.

\textbf{misattribute:} The article misattributes the quotation to a twentieth-century writer.

\textbf{attribute:} The article correctly attributes the quotation to an eighteenth-century philosopher.

\section{Etymology}
\sectionrule
\textbf{Origin:} Misattribute is formed from the prefix mis-, meaning wrongly or incorrectly, and attribute, meaning to assign something to a cause, source, or creator.

\section{Practice exercises}
\sectionrule
\subsection{Word family choice}
Choose the best member of the word family for each blank.

\begin{enumerate}
  \item The museum corrected the \_\_\_\_\_ of the painting after discovering the artist's signature.
  \item Researchers should be careful not to \_\_\_\_\_ an effect to the wrong cause.
\end{enumerate}

\subsection{Correct or incorrect?}
Decide whether each sentence is correct. Correct the inaccurate ones.

\begin{enumerate}
  \item The quotation is commonly misattributed to Einstein.
  \item Doctors initially misattributed the symptoms from stress.
  \item The article misquoted the statement to a famous politician.
  \item The researchers misattributed the improvement to the treatment, but later evidence showed that another factor was responsible.
  \item The historian wrongly misattributed the poem to Shakespeare.
\end{enumerate}

\subsection{Collocation practice}
Complete each sentence with a natural collocation.

\begin{enumerate}
  \item The quotation is often misattributed \_\_\_\_\_ a famous scientist.
  \item Doctors should avoid misattributing physical \_\_\_\_\_ to anxiety without further examination.
  \item The museum discovered that the painting had been \_\_\_\_\_ to the wrong artist.
  \item The error was a clear case of \_\_\_\_\_ because the quotation came from a completely different author.
\end{enumerate}

\subsection{Complete answer key}
\subsubsection{Word family choice}
\begin{enumerate}
  \item \textbf{misattribution.} The museum corrected the misattribution of the painting after discovering the artist's signature. \emph{Use the noun misattribution for an incorrect identification of a work's creator.}
  \item \textbf{misattribute.} Researchers should be careful not to misattribute an effect to the wrong cause. \emph{Use the base verb misattribute after to.}
\end{enumerate}

\subsubsection{Correct or incorrect?}
\begin{enumerate}
  \item \textbf{Correct} \emph{Be misattributed to is the standard pattern for identifying the wrong source of a quotation.}
  \item \textbf{Incorrect}. Doctors initially misattributed the symptoms to stress. \emph{Use misattribute something to a supposed cause.}
  \item \textbf{Incorrect}. The article misattributed the statement to a famous politician. \emph{Use misattribute when the mistake concerns who said something.}
  \item \textbf{Correct} \emph{The sentence correctly uses misattribute for an incorrect explanation of a result.}
  \item \textbf{Incorrect}. The historian misattributed the poem to Shakespeare. \emph{Misattribute already contains the meaning wrongly attribute, so wrongly is unnecessary.}
\end{enumerate}

\subsubsection{Collocation practice}
\begin{enumerate}
  \item \textbf{to.} The quotation is often misattributed to a famous scientist. \emph{Use misattribute something to a person or source.}
  \item \textbf{symptoms.} Doctors should avoid misattributing physical symptoms to anxiety without further examination. \emph{Misattribute symptoms to is a common pattern when discussing an incorrect diagnosis or explanation.}
  \item \textbf{misattributed.} The museum discovered that the painting had been misattributed to the wrong artist. \emph{Be misattributed to means be incorrectly connected with a particular creator.}
  \item \textbf{misattribution.} The error was a clear case of misattribution because the quotation came from a completely different author. \emph{A case of misattribution is an example of identifying the wrong source.}
\end{enumerate}

\vfill
\begin{center}
\small\color{Muted} Created and compiled by \textbf{Amir Kooshky}\\
\href{https://t.me/KooshkyTOEFL}{Telegram Channel}\quad\textbullet\quad
\href{https://www.instagram.com/kooshkytoefl}{Instagram}\quad\textbullet\quad
\href{https://t.me/KooshkyTOEFL\_pv}{Personal Telegram}
\end{center}
\end{document}
